% Section 3 — Augustine Problem + CIT (compressed)

\section{The Augustine Problem and Its Impossibility}
\label{sec:augustine-problem}

\begin{definition}[Augustine Problem]
\label{def:augustine-problem}
Agent $A$ exhibits the Augustine Problem on task family $\tasks$ if
$\mathbb{E}_{\trajectory \sim \pi_A}\!\left[|\twall - \tstep\langle\Delta t\rangle| + |\twall - \tself|\right] > \cce_0$
for some tolerance $\cce_0$. The three projections drift independently across the trajectory distribution.
\end{definition}

The question is not whether agents exhibit it (they do, \S\ref{sec:l2}--\ref{sec:l3}) but whether any conceivable foundation-model training procedure could in principle close it. We argue: not by scale, not by reasoning-compute expansion, not by any post-training intervention on token sequences alone.

\paragraph{Assumption (token-only loss support).}
\label{ass:token-only}
A training procedure is in the \emph{token-only regime} if
$\loss(\theta) = \mathbb{E}_{x \sim \mathcal{D}}[\ell(f_\theta(x_{<i}), x_i)]$
satisfies (i) $\mathcal{D}$ is a distribution over finite token sequences with no field carrying wall-clock duration, and (ii) $\ell$ is a function of the output token distribution alone. Next-token cross-entropy, SFT, RLHF, RL with verifiable rewards, and reasoning supervision all satisfy this.

\begin{theorem}[Chronoception Impossibility, CIT]
\label{thm:cit}
Let $\loss$ be a token-only loss. For any strictly increasing $\phi: \mathbb{R}_{>0} \to \mathbb{R}_{>0}$, $\loss(\theta)$ is invariant under $\twall \mapsto \phi(\twall)$. Equivalently, $\nabla_\theta \loss$ carries zero gradient signal aligning $\twall$ with any internal representation.
\end{theorem}

\begin{proof}[Sketch]
$\nabla_\theta\loss = \mathbb{E}[\partial\ell/\partial f_\theta \cdot \nabla_\theta f_\theta]$. Neither factor depends on the wall-clock time at which the sample was generated: tokens carry no timestamps, $\ell$ is a function of the output distribution alone. The gradient is therefore invariant to any reparametrisation of $\twall$. A policy from this gradient cannot distinguish one wall-clock duration from another.
\end{proof}

\paragraph{Scope.} CIT does not claim there is \emph{no} internal representation of time; it claims any such representation is not produced \emph{by the loss}. Explicitly out of scope: harness-installed chronoception (\S\ref{sec:injection-tell}, a CIT workaround, not refutation), tool-output-installed chronoception (calling \texttt{date}), and future losses with $\twall$ as an explicit argument (\textsc{ChronoStack}). CIT precludes only the inference --- widely held, rarely defended --- that additional scale, longer reasoning, or larger context will close the gap.

\paragraph{Empirical predictions.} \textbf{(P-CIT-1)} $\CAR$ stays bounded away from 1 regardless of capability (\S\ref{sec:l2}). \textbf{(P-CIT-2)} Prompt-level wall-clock information cannot move action-axis behaviour (Prop.~\ref{prop:injection-not-action}). \textbf{(P-CIT-3)} Calibration tooling cannot close self-duration calibration (\S\ref{ssec:calibration}). All three are confirmed by the panel.

\paragraph{Augustine Threshold.} We adopt $\ccestar = 0.20$ as an operational reporting threshold, and we are explicit about its status. It is \emph{not} derived from a measured human baseline. Human duration perception is itself imperfect and strongly scale-dependent \citep{wittmann2009chronoception,eagleman2008time}, which is the qualitative motivation for setting the bar below perfection rather than at it --- but that literature reports \emph{discrimination} thresholds (Weber fractions on whether two intervals differ), not the \emph{retrospective estimation} error that $\confab$ measures, and the two are not interconvertible without an argument we cannot currently supply. We therefore fix $\ccestar = 0.20$ by stipulation: the round figure roughly one order of magnitude below a $2\times$ misestimate ($|\confab| = \log_{10} 2 \approx 0.30$). It was chosen before the panel was run and held fixed afterwards.

We flag this because the threshold is the one place a reader might expect a derivation and not find one. Nothing in the argument depends on the constant. The closest agent sits at $\cce = 0.276$, and the $\tstep$-axis term alone exceeds $0.30$ on every agent in the panel, so every qualitative conclusion here survives any choice of threshold in $[0.1, 0.3]$. A task-matched human panel would let $\ccestar$ be earned rather than stipulated, and is the highest-priority follow-up (\S\ref{sec:limitations}). No panel agent crosses it (\S\ref{sec:epsilon}).

CIT says the gap is unavoidable. To measure ``unavoidable'' we need the instrument next.
